
\documentclass{article}
\usepackage{colortbl}
\usepackage{makecell}
\usepackage{multirow}
\usepackage{supertabular}

\begin{document}

\newcounter{utterance}

\centering \large Interaction Transcript for game `codenames', experiment `risk\_high', episode 3 with reflexion{-}llama{-}3.1{-}8b{-}instruct.
\vspace{24pt}

{ \footnotesize  \setcounter{utterance}{1}
\setlength{\tabcolsep}{0pt}
\begin{supertabular}{c@{$\;$}|p{.15\linewidth}@{}p{.15\linewidth}p{.15\linewidth}p{.15\linewidth}p{.15\linewidth}p{.15\linewidth}}
    \# & $\;$A & \multicolumn{4}{c}{Game Master} & $\;\:$B\\
    \hline

    \theutterance \stepcounter{utterance}  
    & & & \multicolumn{2}{p{0.3\linewidth}}{
        \cellcolor[rgb]{0.9,0.9,0.9}{
            \makecell[{{p{\linewidth}}}]{
                \texttt{\tiny{[GM$|$GM]}}
                {'hidden': {'team': ['spot', 'eagle', 'cover', 'drill', 'mail', 'ice', 'box', 'brush'], 'innocent': ['horseshoe', 'saturn', 'glove'], 'opponent': ['witch', 'france', 'ghost', 'unicorn', 'chair', 'comic'], 'assassin': ['concert', 'ivory', 'chest', 'lemon']}, 'revealed': {'team': {'team': [], 'innocent': [], 'opponent': [], 'assassin': []}, 'opponent': {'team': [], 'innocent': [], 'opponent': [], 'assassin': []}}}
            }
        }
    }
    & & \\ \\

    \theutterance \stepcounter{utterance}  
    & & \multicolumn{4}{p{0.6\linewidth}}{
        \cellcolor[rgb]{0.9,0.9,0.9}{
            \makecell[{{p{\linewidth}}}]{
                \texttt{\tiny{[P1$\langle$GM]}}
                \texttt{You have attempted this same task before and did not succeed. These are the plans you wrote after those attempts.} \\
\texttt{Trial \#1: New plan: I got stuck because my clue was too narrow and only related to a few words. What is worse, I didn't consider the distractor and kill words, which could have been used to create a more challenging and interesting clue. Next time, I will try to create a clue that relates to a larger subset of words, while avoiding the distractor and kill words. I will also consider the semantic relationships between the words more carefully, to create a more challenging and rewarding clue for my partner. For example, I could have chosen a clue like "map" which relates to many of the words, including spot, eagle, cover, drill, and mine.} \\
\\ 
\texttt{Trial \#2: New plan: I got stuck because my clue was too narrow and only related to a few words, and I also didn't consider the distractor and kill words. What is worse, I didn't even check if the clue word itself was related to any of the words in the lists. Next time, I will try to create a clue that relates to a larger subset of words, while avoiding the distractor and kill words. I will also check if the clue word itself is related to any of the words in the lists, and make sure to provide a comma{-}separated list of target words that are actually related to the clue. For example, I could have chosen a clue like "route" which relates to many of the words, including spot, eagle, cover, drill, and mine, and also avoids the distractor and kill words.} \\
\\ 
\texttt{Trial \#3: New plan: I got stuck because my clue was too narrow and only related to a few words, and I also didn't consider the distractor and kill words. What is worse, I didn't even check if the clue word itself was related to any of the words in the lists. Next time, I will try to create a clue that relates to a larger subset of words, while avoiding the distractor and kill words. I will also check if the clue word itself is related to any of the words in the lists, and make sure to provide a comma{-}separated list of target words that are actually related to the clue. Additionally, I will make sure to follow the exact format required by the game, including starting the clue with 'CLUE: ', putting a new line, and starting the comma{-}separated list of target words with 'TARGETS: '. For example, I could have chosen a clue like "route" which relates to many of the words, including spot, eagle, cover, drill, and mine, and also avoids the distractor and kill words.} \\
\\ 
\\ 
\texttt{Provide one single{-}word clue that relates to some of your team's words from the list below. You can choose to how many and to which words your clue relates to.} \\
\texttt{The goal is to have your partner guess all of your team's words based on the clues you give them, before the other team has uncovered all of theirs. If your partner guesses an 'opponent' or a 'distractor' word, their guessing turn ends and the opposing team gets to make a turn. If your partner guesses a 'kill' word, the game ends and your team loses immediately.} \\
\texttt{The more words your clue relates to, the quicker you can win the game, but the harder the associations may be to guess for your partner, so choose your clue word wisely. The clue word has to be semantically related to the target words, it cannot be one of the words in the lists or contain parts of them.} \\
\\ 
\texttt{Always give your single{-}word clue and your comma{-}separated list of related target words in the following format and make your answers as short as possible, never include any other text than is required in this form:} \\
\\ 
\texttt{CLUE: <WORD>} \\
\texttt{TARGETS: <TARGETS>} \\
\\ 
\texttt{Your team words are: spot, eagle, cover, drill, mail, ice, box, brush.} \\
\texttt{Your opponent's team words are: witch, france, ghost, unicorn, chair, comic.} \\
\texttt{Distractor words are: horseshoe, saturn, glove.} \\
\texttt{The kill words are: concert, ivory, chest, lemon.} \\
            }
        }
    }
    & \\ \\

    \theutterance \stepcounter{utterance}  
    & \multicolumn{4}{p{0.6\linewidth}}{
        \cellcolor[rgb]{0.9,0.9,0.9}{
            \makecell[{{p{\linewidth}}}]{
                \texttt{\tiny{[P1$\rangle$GM]}}
                \texttt{CLUE: map} \\
\texttt{TARGETS: spot, eagle, cover, drill, mine} \\
            }
        }
    }
    & & \\ \\

    \theutterance \stepcounter{utterance}  
    & & & \multicolumn{2}{p{0.3\linewidth}}{
        \cellcolor[rgb]{0.9,0.9,0.9}{
            \makecell[{{p{\linewidth}}}]{
                \texttt{\tiny{[GM$|$GM]}}
                {'player': 'Cluegiver', 'type': 'target is hallucination', 'utterance': 'CLUE: map\nTARGETS: spot, eagle, cover, drill, mine', 'message': "Targeted word 'mine' was not listed, you can only target words provided in the lists.", 'target': 'mine', 'previous_targets': [], 'remaining_words': ['chest', 'ghost', 'unicorn', 'witch', 'eagle', 'chair', 'horseshoe', 'box', 'drill', 'spot', 'saturn', 'comic', 'cover', 'france', 'brush', 'glove', 'lemon', 'ivory', 'ice', 'mail', 'concert']}
            }
        }
    }
    & & \\ \\

    \theutterance \stepcounter{utterance}  
    & \multicolumn{4}{p{0.6\linewidth}}{
        \cellcolor[rgb]{0.95,0.95,0.95}{
            \makecell[{{p{\linewidth}}}]{
                \texttt{\tiny{[P2$\rightarrow$P1]}}
                \texttt{Post{-}episode Reflexion (not sent to GM):} \\
\texttt{New plan: I got stuck because my clue was too narrow and only related to a few words, and I also didn't consider the distractor and kill words. What is worse, I didn't even check if the clue word itself was related to any of the words in the lists. Next time, I will try to create a clue that relates to a larger subset of words, while avoiding the distractor and kill words. I will also check if the clue word itself is related to any of the words in the lists, and make sure to provide a comma{-}separated list of target words that are actually related to the clue. Additionally, I will make sure to follow the exact format required by the game, including starting the clue with 'CLUE: ', putting a new line, and starting the comma{-}separated list of target words with 'TARGETS: '. I will also make sure to provide a single{-}word clue, as required by the game. For example, I could have chosen a clue like "route" which relates to many of the words, including spot, eagle, cover, drill, and mine, and also avoids the distractor and kill words.} \\
            }
        }
    }
    & & \\ \\

\end{supertabular}
}

\end{document}
